\section{Requisitos Implementados}

% FONTES: docs/ARQUITETURA.md §5 (5.1 gRPC, 5.2 Bully, 5.3 Lamport, 5.4 Ricart-Agrawala),
%         docs/algoritmos/*.md (algoritmo + diagrama de sequencia de cada um),
%         docs/classes/*.md (detalhe de implementacao por classe)
% O enunciado pede NO MINIMO DOIS requisitos; aqui sao quatro. Cada subsecao deve caber em
% ~1/3 de pagina: algoritmo, mensagens e UMA decisao de projeto.

\subsection{RPC -- gRPC}
% src/main/proto/node.proto, rpc/PeerNetwork.java, rpc/NodeServiceImpl.java, rpc/ClusterAuth.java
\pendente{servi\c{c}o \'unico; deadlines; token no interceptor; 64 MiB}

\subsection{Elei\c{c}\~ao de l\'ider -- Bully}
% election/BullyElection.java, election/LeaderMonitor.java, docs/algoritmos/bully.md
\pendente{3 passos; maior id; 5 heartbeats; contesta\c{c}\~ao de id menor; failover (Tabela 2)}

\subsection{Rel\'ogio l\'ogico -- Lamport}
% clock/LamportClock.java, docs/algoritmos/lamport.md
\pendente{tick no envio, max+1 no recebimento; ordena pedidos e numera logs}

\subsection{Exclus\~ao m\'utua -- Ricart-Agrawala}
% mutex/RicartAgrawala.java, docs/algoritmos/ricart-agrawala.md
\pendente{RELEASED/WANTED/HELD; prioridade (ts, id); REPLY carrega request\_ts; 2(N-1) mensagens;
falha = permiss\~ao (risco assumido, vira desafio na se\c{c}\~ao 5)}

\begin{figure}[t]
  \centering
  \scalebox{0.85}{% Figura 2 -- Ricart-Agrawala com tres participantes.
% Uso: % Figura 2 -- Ricart-Agrawala com tres participantes.
% Uso: % Figura 2 -- Ricart-Agrawala com tres participantes.
% Uso: \input{figuras/sequencia-ra} dentro de um ambiente figure.
\begin{tikzpicture}[
    font=\footnotesize,
    cab/.style={draw, rounded corners=2pt, minimum width=1.15cm, minimum height=0.6cm, fill=black!6},
    seta/.style={-{Latex[length=1.7mm]}, semithick},
    rot/.style={font=\scriptsize, inner sep=1.5pt, fill=white}
  ]

  % Cabecalhos e linhas de vida
  \node[cab] (h1) at (0,0)   {N1};
  \node[cab] (h2) at (3.2,0) {N2};
  \node[cab] (h3) at (6.4,0) {N3};
  \foreach \x in {0,3.2,6.4} { \draw[dashed, gray] (\x,-0.32) -- (\x,-4.45); }

  % REQUEST do N1 para os dois peers
  \draw[seta] (0,-1.00) -- node[rot, above] {REQUEST $(ts=5)$} (3.2,-1.00);
  \draw[seta] (0,-1.55) -- node[rot, above] {REQUEST $(ts=5)$} (6.4,-1.55);

  % Permissoes
  \draw[seta] (3.2,-2.10) -- node[rot, above] {REPLY} (0,-2.10);
  \draw[seta] (6.4,-2.55) -- node[rot, above] {REPLY} (0,-2.55);

  % N1 entra na secao critica com as duas permissoes
  \fill[black!12] (-0.42,-3.00) rectangle (0.42,-4.15);
  \node[font=\scriptsize, align=center, text width=1.05cm] at (0,-3.58)
    {se\c{c}\~ao\\cr\'itica};

\end{tikzpicture}
 dentro de um ambiente figure.
\begin{tikzpicture}[
    font=\footnotesize,
    cab/.style={draw, rounded corners=2pt, minimum width=1.15cm, minimum height=0.6cm, fill=black!6},
    seta/.style={-{Latex[length=1.7mm]}, semithick},
    rot/.style={font=\scriptsize, inner sep=1.5pt, fill=white}
  ]

  % Cabecalhos e linhas de vida
  \node[cab] (h1) at (0,0)   {N1};
  \node[cab] (h2) at (3.2,0) {N2};
  \node[cab] (h3) at (6.4,0) {N3};
  \foreach \x in {0,3.2,6.4} { \draw[dashed, gray] (\x,-0.32) -- (\x,-4.45); }

  % REQUEST do N1 para os dois peers
  \draw[seta] (0,-1.00) -- node[rot, above] {REQUEST $(ts=5)$} (3.2,-1.00);
  \draw[seta] (0,-1.55) -- node[rot, above] {REQUEST $(ts=5)$} (6.4,-1.55);

  % Permissoes
  \draw[seta] (3.2,-2.10) -- node[rot, above] {REPLY} (0,-2.10);
  \draw[seta] (6.4,-2.55) -- node[rot, above] {REPLY} (0,-2.55);

  % N1 entra na secao critica com as duas permissoes
  \fill[black!12] (-0.42,-3.00) rectangle (0.42,-4.15);
  \node[font=\scriptsize, align=center, text width=1.05cm] at (0,-3.58)
    {se\c{c}\~ao\\cr\'itica};

\end{tikzpicture}
 dentro de um ambiente figure.
\begin{tikzpicture}[
    font=\footnotesize,
    cab/.style={draw, rounded corners=2pt, minimum width=1.15cm, minimum height=0.6cm, fill=black!6},
    seta/.style={-{Latex[length=1.7mm]}, semithick},
    rot/.style={font=\scriptsize, inner sep=1.5pt, fill=white}
  ]

  % Cabecalhos e linhas de vida
  \node[cab] (h1) at (0,0)   {N1};
  \node[cab] (h2) at (3.2,0) {N2};
  \node[cab] (h3) at (6.4,0) {N3};
  \foreach \x in {0,3.2,6.4} { \draw[dashed, gray] (\x,-0.32) -- (\x,-4.45); }

  % REQUEST do N1 para os dois peers
  \draw[seta] (0,-1.00) -- node[rot, above] {REQUEST $(ts=5)$} (3.2,-1.00);
  \draw[seta] (0,-1.55) -- node[rot, above] {REQUEST $(ts=5)$} (6.4,-1.55);

  % Permissoes
  \draw[seta] (3.2,-2.10) -- node[rot, above] {REPLY} (0,-2.10);
  \draw[seta] (6.4,-2.55) -- node[rot, above] {REPLY} (0,-2.55);

  % N1 entra na secao critica com as duas permissoes
  \fill[black!12] (-0.42,-3.00) rectangle (0.42,-4.15);
  \node[font=\scriptsize, align=center, text width=1.05cm] at (0,-3.58)
    {se\c{c}\~ao\\cr\'itica};

\end{tikzpicture}
}
  \caption{Ricart-Agrawala com tr\^es participantes: N1 s\'o entra na se\c{c}\~ao cr\'itica depois
  de receber as duas permiss\~oes. O par $(ts, id)$ decide quem espera: um pedido mais novo pode
  ser adiado e respondido s\'o na sa\'ida da se\c{c}\~ao cr\'itica.}
  \label{fig:ra}
\end{figure}
